\documentclass[a4paper,12pt]{article}
\usepackage[utf8]{inputenc}
\usepackage[T1]{fontenc}
\usepackage{amsmath, amssymb}
\usepackage{geometry}
\geometry{margin=2cm}
\usepackage{xcolor}
\usepackage{tcolorbox}
\tcbuselibrary{skins, breakable}
\usepackage[colorlinks,citecolor=blue]{hyperref}



% ------------------------------------------------------
% تعریف محیط نظر داور
% ------------------------------------------------------
\newtcolorbox{reviewerscomment}{
    colback=blue!5,          % رنگ زمینه بسیار کم‌رنگ آبی
    colframe=blue!30!yellow,  % رنگ کادر
    arc=3mm,                 % گردی گوشه‌ها
    left=5mm, right=5mm, top=3mm, bottom=3mm,
    fonttitle=\bfseries,
    title=نظر داور,
    breakable                % امکان شکستن در چند صفحه
}

% ------------------------------------------------------
% تعریف محیط پاسخ شما
% ------------------------------------------------------
\newtcolorbox{myresponse}{
    colback=green!5,           % رنگ زمینه بسیار کم‌رنگ قرمز
    colframe=green!70!black,   % 
    arc=3mm,
    left=5mm, right=5mm, top=3mm, bottom=3mm,
    fonttitle=\bfseries,
    title=پاسخ به داور,
    breakable
}

% برای شماره‌گذاری خودکار نظرات داوران
\newcounter{commentcounter}
\setcounter{commentcounter}{0}
\newcommand{\reviewercomment}[1]{%
    \stepcounter{commentcounter}%
    \begin{reviewerscomment}
        \textbf{نظر \arabic{commentcounter}:} #1
    \end{reviewerscomment}
}

% ------------------------------------------------------
% شروع سند
% ------------------------------------------------------
\begin{document}

\title{پاسخ به نظرات داوران مقاله\\
    \large
    ''از فیلترهای کلاسیک تا شبکه‌های عصبی پیچشی: مبانی ریاضی و تکامل انتقال سَبْک در تصاویر``}
\author{محمود امین‌طوسی}
\date{\today}
\maketitle

%\section*{مقدمه}
در ابتدا از داور محترم به خاطر زمانی که صرف بررسی دقیق این نوشتار کرده‌اند، صمیمانه سپاسگزاری می‌کنم. 
 داور گرامی چنین مرقوم فرموده بودند:
\begin{quotation}\bf
این یک مقاله مروری است که روند و مبانی انتقال سبک در تصاویر را تشریح می کند. مقاله بخوبی و زیبایی تدوین شده است و اشکالات نگارشی کمی دارد. بعضی از این اشکالات در فایلpdf  مقاله برجسته شده است و به انضمام این متن ارسال می گردد. در ضمن تکرار معادل انگلیسی برخی زیر نویس ها چک شود. بطور مثال معادل یادگیری عمیق بصورت زیرنویس در صفحات ۱ و ۷ تکرار شده است.
سوالات و پیشنهادات زیر برای بهبود مقاله ارائه می گردد.
\end{quotation}


ضمن سپاس از نظر لطف ایشان و دیدن جنبه‌های مثبت کار،  در ادامه، نظرات ایشان همراه با پاسخ‌ها و اقدامات انجام‌شده ارائه شده است. ضمناً مواردی که در فایل pdf‌ ارسالی برجسته فرموده بودند در نسخه بازبینی شده اعمال شده است.

% ------------------------------------------------------
% نظر داور 1
% ------------------------------------------------------



\reviewercomment{
مقاله با وجود عنوان گسترده، فاقد یک پرسش پژوهشی مشخص است. آیا هدف مقاله مرور تاریخی مفاهیم ریاضی مرتبط با انتقال سبک در تصاویر است یا ارائه ی دیدگاهی تحلیلی درباره ی تکامل انتقال سبک عصبی؟ لطفا هدف اصلی مقاله را به وضوح در مقدمه بیان کنید و ساختار مقاله را بر اساس آن تنظیم نمایید.}

\begin{myresponse}
با تشکر از توجه داور محترم؛ مقاله حاضر از نوع «علمی-ترویجی» است و برای مجله «ریاضی و جامعه» که رسالت آن عمومی‌سازی ریاضیات و نشان دادن کاربردهای جذاب آن در زندگی و فناوری روز است، نگاشته شده است. هدف اصلی من این بوده است که بر پایه دانش ریاضیات دوره کارشناسی، خوانندگانی که با ریاضیات کلاسیک آشنا هستند را از طریق یک مثال جذاب با مفاهیم یادگیری ماشین آشنا کنم؛‌ و به نحوی محدودیت‌های ریاضیات سنتی در مواجهه با مسائل جدید داده محور را بیان کنم.
یعنی نه می خواستم به پرسش پژوهشی مشخصی پاسخ دهم و نه تحلیل جامعی درباره انتقال سبک عصبی داشته باشم. سعی من در عنوان و متن نوشتار این بوده است که خواننده‌ای که دانشی از ریاضیات دارد را جذب کند و خواننده را از دنیای قدیم به دنیای جدید بکشاند.

البته تصور من این هست که داور محترم نگاه «پژوهشی» به مقاله داشته‌اند،‌در حالیکه رویکرد آن «علمی ترویجی» بوده است.
شاید نوشتار من به قدر کافی گویا نبوده است، لذا بخش پایانی مقدمه با هدف بیان بهتر هدف مقاله، بازنویسی شد.
اگر توضیح خاصی لازم هست به مقاله اضافه شود یا عنوان اصلاح شود، از پیشنهادات سازنده داور گرامی استقبال می‌کنم.
\end{myresponse}

% ------------------------------------------------------
% نظر داور 2
% -----------------------------------------------------
\reviewercomment{
در بخش ۴، ماتریس گرام را به عنوان زبان ریاضی سبک معرفی می کند، اما به محدودیت های آن اشاره نمی کند. برای مثال، آیا ماتریس گرام می تواند تفاوت میان سبک های مشابه را به خوبی تشخیص دهد؟ چه جایگزین هایی برای آن پیشنهاد شده و مزایا و معایب هر یک چیست؟}
\begin{myresponse}

با سپاس از نقد داور محترم؛ باید یادآوری کنم که تم اصلی این نوشتار، همان‌گونه که در بالاتر بیان شد، «علمی-ترویجی» است و تمرکز آن بر برقراری پلی میان دانش ریاضی کلاسیک و فناوری مدرن یادگیری عمیق می‌باشد، نه ارائهٔ یک تحلیل جامع و نقادانه از روش‌های انتقال سبک. با این حال، مطابق نظر ایشان، یک پاراگراف در این خصوص به اواخر بخش چهار اضافه شد.

\end{myresponse}

\bibliographystyle{vancouver} 
\bibliography{refs}



\end{document}